\section{A free excitation direction: limit, warning, and PSM slice}
\label{sec:loss-degeneracy}

Write \(a=3R_s/2>0\) and \(b=R_{\mathrm{exc}}>0\).  At loss level
\(P_{\mathrm{Cu},\max}\), the three semiaxes are
\begin{equation}
 I_d=I_q=\sqrt{\frac{P_{\mathrm{Cu},\max}}{a}},\qquad
 I_e=\sqrt{\frac{P_{\mathrm{Cu},\max}}{b}}.
 \label{eq:loss-semiaxes}
\end{equation}
Consequently \(b\downarrow0\) sends \(I_e\) to infinity while leaving the
stator radius fixed.  The positive-definite loss ellipsoid converges on every
bounded window to the cylinder
\begin{equation}
 \boxed{a(i_d^2+i_q^2)=P_{\mathrm{Cu},\max},
 \qquad i_{\mathrm{exc}}\text{ free}.}
 \label{eq:loss-cylinder}
\end{equation}
The matrix \(\matr R=\operatorname{diag}(a,a,b)\) loses rank and
\(\matr R^{-1/2}\) ceases to exist.  The divergence is therefore visible both
before and after whitening: a torque-producing direction has lost its price.

\begin{figure}[tb]
\centering
\begin{tikzpicture}[scale=.92,>=Latex]
 \foreach \x/\h/\lab in {0/1.15/{b>0},4.1/1.85/{b\downarrow0},8.2/2.35/{b=0}}{
   \draw[thin,->] (\x-1.25,0)--(\x+1.35,0) node[right] {$i_d$};
   \draw[thin,->] (\x,-1.25)--(\x,2.45) node[above] {$i_{\mathrm{exc}}$};
   \ifdim \x pt<8pt
     \draw[blue,thick,fill=blue!8] (\x,0.55) ellipse (1.05 and \h);
     \draw[blue!55,dashed] (\x-1.05,.55)--(\x+1.05,.55);
   \else
     \draw[blue,thick] (\x-1.05,-1.05)--(\x-1.05,2.15);
     \draw[blue,thick] (\x+1.05,-1.05)--(\x+1.05,2.15);
     \draw[blue!55,dashed] (\x-1.05,.55)--(\x+1.05,.55);
     \draw[blue,->] (\x-1.05,2.0)--(\x-1.05,2.4);
     \draw[blue,->] (\x+1.05,2.0)--(\x+1.05,2.4);
     \draw[blue,->] (\x-1.05,-.9)--(\x-1.05,-1.3);
     \draw[blue,->] (\x+1.05,-.9)--(\x+1.05,-1.3);
   \fi
   \node[font=\small] at (\x,-1.65) {$\lab$};
 }
 \draw[->,very thick,orange] (1.5,.55)--(2.6,.55)
   node[midway,above,font=\scriptsize] {$I_e\propto b^{-1/2}$};
 \draw[->,very thick,orange] (5.6,.55)--(6.7,.55);
\end{tikzpicture}
\caption{Loss-geometry degeneracy in an \(i_d\)-\(i_{\mathrm{exc}}\)
section.  The full three-dimensional bodies are ellipsoids of revolution about
the excitation axis; as \(R_{\mathrm{exc}}=b\) decreases, the excitation
semiaxis diverges and the limit is the circular cylinder
\cref{eq:loss-cylinder}.}
\label{fig:loss-degeneracy}
\end{figure}


This is not automatically a PSM limit.  If excitation remains controllable,
the loss-only model admits \(i_q=\varepsilon\),
\(i_{\mathrm{exc}}=c/\varepsilon\), and \(i_d=0\).  Then
\begin{equation}
 M=k_TL_{\mathrm{exc}}c+O(\varepsilon),\qquad
 P_{\mathrm{Cu}}=a\varepsilon^2,
 \label{eq:free-field-scaling}
\end{equation}
so a finite excitation torque can approach zero modeled copper loss.  With a
fixed nonzero \(i_q\), torque is unbounded as \(i_{\mathrm{exc}}\to\infty\).
At nonzero speed, however,
\(u_q=\omega_e(L_di_d+L_{\mathrm{exc}}i_{\mathrm{exc}})+R_si_q\)
usually blocks this sequence: canceling the growing field term requires an
equally growing \(i_d\), which is still priced by \(a>0\).  At exactly zero
speed the stator voltage does not see excitation, so the degeneracy survives.
An explicit field-current or field-voltage limit, magnetic saturation, rotor
thermal constraints, or a finite excitation converter also regularizes it.

The PSM-like limit instead fixes a slice,
\begin{equation}
 i_{\mathrm{exc}}=i_{e0},\qquad
 \psi_{\mathrm{pm}}=L_{\mathrm{exc}}i_{e0}=\text{constant}.
 \label{eq:fixed-excitation-limit}
\end{equation}
On that plane, the EESM torque becomes
\(M=k_Ti_q(\Delta L i_d+\psi_{\mathrm{pm}})\), exactly the affine PSM
law.  If \(b\to0\) as well, the constant offset \(b i_{e0}^2\) disappears
from loss while its flux effect remains.  Fixing excitation removes a degree
of freedom; setting its resistance to zero removes a cost.  These operations
do not commute conceptually.

\begin{figure}[tb]
\centering
\begin{tikzpicture}[scale=.95,>=Latex]
 \begin{scope}[xshift=-3.3cm]
  \draw[thin,->] (-2,0)--(2.2,0) node[right] {$i_d$};
  \draw[thin,->] (0,-1.55)--(0,2.15) node[above] {$i_{\mathrm{exc}}$};
  \draw[thin,->] (-.9,-.7)--(1.25,.95) node[above right] {$i_q$};
  \draw[blue,thick,fill=blue!7] (0,.25) ellipse (1.7 and 1.25);
  \draw[orange,thick,fill=orange!12,opacity=.7] (-1.75,.75)--(1.75,.75)--(1.25,1.15)--(-2.25,1.15)--cycle;
  \draw[red,very thick] (-1.42,.94) arc[start angle=185,end angle=355,x radius=1.45,y radius=.42];
  \node[orange!80!black,font=\scriptsize] at (1.05,1.45) {$i_{\mathrm{exc}}=i_{e0}$};
  \node[font=\scriptsize] at (0,-1.9) {EESM current space};
 \end{scope}
 \draw[->,very thick] (-.55,.3)--(.65,.3) node[midway,above,font=\scriptsize] {slice};
 \begin{scope}[xshift=3.2cm]
  \draw[thin,->] (-2,0)--(2.2,0) node[right] {$i_d$};
  \draw[thin,->] (0,-1.55)--(0,1.75) node[above] {$i_q$};
  \draw[blue,thick,fill=blue!7] (0,0) circle (1.25);
  \draw[red,very thick] plot[smooth,domain=-1.7:1.7,samples=80]
    ({\x},{.52/(\x+2.05)});
  \node[font=\scriptsize,align=center] at (0,-1.9)
   {PSM plane\\$\psi_{\mathrm{pm}}=L_{\mathrm{exc}}i_{e0}$};
 \end{scope}
\end{tikzpicture}
\caption{A PSM-like model is a fixed-excitation plane through the EESM current
space.  The plane turns controllable excitation into the fixed affine flux
offset \(\psi_{\mathrm{pm}}\); it is not produced by merely setting
\(R_{\mathrm{exc}}\) to zero.}
\label{fig:eesm-psm-slice}
\end{figure}


\paragraph{Limits as model diagnostics.}
Whenever a capability diverges, the first question should be: which
physically effective direction has just become free?  The same rule diagnoses
other limits.  As \(R_s\to0\), stator current loses its thermal price and
whitening fails in two directions.  As \(L_{\mathrm{exc}}\to0\), field
current remains costly but stops producing torque.  As \(\Delta L\to0\), the
reluctance mechanism vanishes.  Letting \(P_{\mathrm{Cu},\max}\to\infty\) or
\(U_{\max}\to\infty\) removes a bounding surface rather than improving a
physical machine.  A finite answer after such a removal must come from a
remaining constraint, not from the optimizer's search box.

